\section{Exact local algebras and extraction limits}\label{sec:exact-law}
This appendix separates two refinements from the approximate-relation theorem. The first assumes exact local relations and has a sharp leading coefficient. The second limits a full-state self-testing modulus. Neither refinement is needed for Theorem~D.

\subsection{The sharp exact-algebra law}
\begin{theorem}[Exact local relations; three mirrors]\label{thm:exact-law}
Suppose both parties obey \eqref{eq:R1}--\eqref{eq:R3} as operator identities. On a normal purification $\Omega$, set
\[
 \Delta_3=\|(A_L^*-A_R/2)\Omega\|^2
 +\|(B_L^*-B_R)\Omega\|^2+\|(C_L-C_R)\Omega\|^2.
\]
There is an absolute constant $C_0$ such that, for sufficiently small $\Delta_3>0$ and either party $s$,
\[
 S(\rho_s)\ge\frac23\log_2(1/\Delta_3)
       -8\log_2\log_2(1/\Delta_3)-C_0.
\]
In particular the same bound holds with $\Delta_3$ replaced by the sum $\Delta$ of all five squared mirror defects. For that five-mirror convention the leading coefficient $2/3$ is sharp, uniformly over representations and normal states.
\end{theorem}
\begin{proof}
Put $\delta=\sqrt{\Delta_3}$ and $\ell=\log_2(1/\delta)$. The spectral argument in the first paragraph of the proof of Lemma~\ref{lem:normal} uses only the local identities: $P,Q$ are orthogonal reducing spectral subprojections, $B$ is a partial isometry, and $AA^*,BB^*,P$ are an orthogonal partition. An interval in $(1/2,1)$ would be transported isometrically by $B$ above one. The same reasoning on $1_{\{1\}}(C)-P$ gives $E_0=P$; on $1_{\{1/2\}}(C)-Q$ it gives $E_1=Q$, since the image would lie in $P$ and in the orthogonal range of $B$. Repeated upward transport excludes every interval $(2^{-h-1},2^{-h})$. Thus
\[
 I=E_\infty+\sum_{h\ge0}E_h,\qquad
 E_h=1_{\{2^{-h}\}}(C),\quad E_\infty=1_{\{0\}}(C).
\]
This step makes no assertion that the zero-counter mass vanishes.

For $k\ge0$ put $F_k=1_{[0,2^{-k}]}(C)$ and $t_k=\|F_k^L\Omega\|^2$, so $t_0=1$. For $k\ge1$, the range partition and covariance give the exact identity
\[
 t_k=\|F_{k-1}^LA_L^*\Omega\|^2
           +\|F_{k+1}^LB_L^*\Omega\|^2.
\]
Replace the first adjoint by $A_R/2$ and the second by $B_R$. Their squared-norm errors are at most $3\delta/2$ and $2\delta$, respectively. Opposite-party commutation and the exact source relations give the first squared norm $t_{k-1}/4$ and the second
\[
 t_{k+1}-\|F_{k+1}^L(P_R+Q_R)\Omega\|^2.
\]
On this last joint spectral support, $C_L\le1/4$ and $C_R\ge1/2$. The commuting spectral-gap estimate bounds its mass by $16\delta^2$. Hence, for $\delta\le1$,
\[
 |t_k-t_{k-1}/4-t_{k+1}|\le20\delta.
\]
The homogeneous solution from $t_0,t_1$ is $[t_0+(2t_1-t_0)k]2^{-k}$, at most $(k+1)2^{-k}$. The impulse kernel has total mass $\sum_{j\ge0}(j+1)2^{-j}=4$. Therefore
\begin{equation}\label{eq:exact-tail}
 t_k\le(k+1)2^{-k}+80\delta.
\end{equation}
Only the three displayed mirrors have been used.

Let $D_\infty$ project onto equal finite heights together with equal zero-counter sectors. If unequal counters have at least one height at most $m$, their distance is at least $2^{-m-1}$. All remaining unequal pairs have left counter at most $2^{-m-1}$. Thus
\[
 \|(I-D_\infty)\Omega\|^2
 \le4^{m+1}\delta^2+(m+2)2^{-m-1}+80\delta.
\]
With $m=\lfloor2\ell/3\rfloor$, this is at most $c_1(\ell+1)\delta^{2/3}$ for an absolute $c_1$. For $D_H=\sum_{h=0}^HE_h^LE_h^R$ and $\Omega_H=D_H\Omega$, \eqref{eq:exact-tail} consequently gives
\begin{equation}\label{eq:exact-match}
 \|\Omega-\Omega_H\|\le
 c_2\sqrt{\ell+1}\,\delta^{1/3}
 +c_2\sqrt{H+2}\,2^{-H/2}.
\end{equation}
The estimate includes zero-counter mass and is independent of a smallest Schmidt value.

Retain $a=A\sum_{h=0}^{H-2}E_h$ and $b=B\sum_{h=2}^HE_h$. Their columns are isometries with orthogonal incoming ranges, and $aa^*+bb^*=\sum_{h=1}^{H-1}E_h$. On matched heights let
\[
 T_H=\tfrac12a_L\otimes a_R+b_L\otimes b_R,
 \qquad \xi_j=E_j^LE_j^R\Omega\quad(j=0,H).
\]
There is no incoming retained edge at either boundary. Define
\[
 \Pi_+=\sum_{h=1}^{H-1}E_{h-1}^LE_h^R,\qquad
 \Pi_-=\sum_{h=1}^{H-1}E_{h+1}^LE_h^R.
\]
Covariance proves, as exact vector identities,
\[
 (a_L^*-a_R/2)\Omega_H=\Pi_+(A_L^*-A_R/2)\Omega,
\quad
 (b_L^*-b_R)\Omega_H=\Pi_-(B_L^*-B_R)\Omega.
\]
Indeed selecting the displayed output heights forces matched input heights in each term. Multiply these identities by $a_L,b_L$ and use the interior row partition. Their ranges are orthogonal, so
\[
 (I-T_H)\Omega_H=\xi_0+\xi_H+r_H,
 \qquad\|r_H\|\le\delta.
\]
In particular the height-matching error has not passed through the inverse. Lemma~\ref{rnd:lem:renewal} gives
\[
 y_H=(I-T_H)^{-1}(\xi_0+\xi_H),\qquad
 \|y_H-\Omega_H\|\le4H^2 2^{H/2}\delta.
\]
Combining with \eqref{eq:exact-match}, and writing the resulting bound as $\zeta_H$, yields
\[
 \zeta_H\le c_3\left(
 \sqrt{\ell+1}\,\delta^{1/3}
 +\sqrt{H+2}\,2^{-H/2}
 +H^2 2^{H/2}\delta\right).
\]
Choose $H=\lfloor2\ell-12\log_2(\ell+2)-c_4\rfloor$. For sufficiently large $\ell$, $H\ge2$ and $\zeta_H(N_{\rm hist}+2)$ is bounded by an absolute constant: use $N_{\rm hist}+2\le18H^3$ and $2^{H/2}\delta\le2^{-c_4/2}(\ell+2)^{-6}$. The first two terms decay exponentially in $\ell$ after multiplication by $H^3$. Also $\zeta_H<1/2$ and $y_H\ne0$ for sufficiently large $\ell$.

The local finite-history isometry and Lemma~\ref{rnd:lem:entropy-comparison}, applied on either side, now give
\[
 S(\rho_s)\ge\tfrac23H-O(1)
 =\tfrac43\ell-8\log_2(\ell+2)-O(1).
\]
The path construction is symmetric, so this gives both marginal bounds even though the renewal residual used only the left adjoint mirrors. Since $\log_2(1/\Delta_3)=2\ell$, the stated remainder follows. If all three defects vanish, apply the same argument with every positive $\delta$ to obtain infinite entropy. Replacing $\Delta_3$ by the larger $\Delta$ changes the lower bound only in the sufficiently small range, where its right side is increasing in $\log_2(1/\Delta)$.

For sharpness use the canonical height-$H$ state on the unchanged exact local representation. $B$ and $C$ mirrors are exact. Each of the two $A$ mirrors has squared norm $q_H/4$, since only an omitted $A$ child contributes, with coefficient $1/2$. Thus $\Delta=q_H/2$. Equations~\eqref{eq:v11-height-formulas} and \eqref{eq:v11-height-entropy}, with the fixed scaffold term removed, give entropy $2H/3+O(1)$ and $\log_2(1/\Delta)=H+\log_2H+O(1)$. No larger uniform leading coefficient can hold.
\end{proof}

\begin{corollary}[The boundary-law class]\label{cor:boundary-law}
For the prefix-closed, $B$-closed canonical restrictions of Section~\ref{sec:v11-height}, the infimum entropy at boundary ratio at most $q$ has leading term $(2/3)\log_2(1/q)$. The same leading coefficient holds for that class of strategies for $B$ and $G$, with their fixed normalizations.
\end{corollary}
\begin{proof}
On the unamplified word state, prefix closure and $B$ closure make both $B$ mirrors exact; $C$ is diagonal and matched. Each $A$ mirror has squared norm $q_S/4$, by the same omitted-child calculation. Hence $\Delta_S=q_S/2$ for every such set, including infinite sets of finite entropy. Theorem~\ref{thm:exact-law} gives the lower leading term uniformly over the class. Height cutoffs give the upper term, and finite height-and-size cutoffs approximate them with strict loss slack. The scaffold adds a fixed entropy; \eqref{eq:v11-boundary} multiplies loss by a fixed positive constant. Neither changes a leading logarithmic coefficient.
\end{proof}

\subsection{A limit on full-state extraction}
\begin{proposition}[No faster full-state self-testing modulus]\label{prop:modulus}
For sufficiently small $\eta>0$, a pure state $\phi$ within vector distance $\eta$ of $\psicat\otimes\xi$, after local isometries and with arbitrary auxiliary vector $\xi$, has marginal entropy at least $c\eta^{-2}$ for an absolute $c>0$. Consequently a uniform vector extraction modulus for all near-perfect strategies of $B$ or $G$ must satisfy
\[
 \eta(\varepsilon)\ge c'\bigl(\log_2(1/\varepsilon)\bigr)^{-1/2}
\]
for sufficiently small $\varepsilon$, with a fixed $c'>0$. In particular no positive H\"older rate is possible.
\end{proposition}
\begin{proof}
Let $p$ be the ordered Schmidt probabilities of $\psicat$, and $d_N=\sum_{n=0}^N\Cat_n$. The elementary Catalan bound $\Cat_n4^{-n}\ge[2(n+1)^{3/2}]^{-1}$ gives
\[
 \tau_p(d_N):=\sum_{i>d_N}p_i\ge\frac1{2\sqrt{N+2}},
 \qquad\log_2d_N\ge N-1.
\]
For the first inequality sum the bound over $n>N$ and compare with the integral of $x^{-3/2}/4$ from $N+2$ to infinity. For the second use $\Cat_N\ge2^{N-1}$, which follows inductively from the two endpoint terms in the Catalan convolution (with $N=1$ checked directly).

Tensoring by an auxiliary Schmidt law cannot decrease an ordered tail. To see this, write $p\otimes q=\sum_jq_j(p\otimes e_j)$. The sum of the largest $d$ entries is convex and is the same for each vector $p\otimes e_j$, so it is at most the sum of the largest $d$ entries of $p$. The conclusion also holds for countably many $q_j$ by monotone approximation.

For any coefficient Hilbert--Schmidt operator, its distance to operators of rank at most $d$ is the square root of its ordered tail. This follows by projecting onto the $d$ largest singular directions for the upper bound, and by the variational characterization of the sum of the largest $d$ squared singular values for the lower bound. Distance to a fixed set is $1$-Lipschitz. Therefore, for the Schmidt law $p_\phi$,
\[
 \sqrt{\tau_{p_\phi}(d)}\ge\sqrt{\tau_p(d)}-\eta.
\]
Take $N=\lfloor(256\eta^4)^{-1}\rfloor$ and $\eta\le1/8$. Then $N+2\le(128\eta^4)^{-1}$, so $\tau_p(d_N)\ge4\eta^2$ and $\tau_{p_\phi}(d_N)\ge\eta^2$. Since the probabilities after position $d_N$ are at most $1/d_N$,
\[
 S(\phi_L)\ge\tau_{p_\phi}(d_N)\log_2d_N
 \ge\frac1{512\eta^2}.
\]
This remains true with an untouched purifying environment in the auxiliary system, since the cut is the original left marginal against everything else.

Apply the estimate to the pure finite height-and-size strategies in Theorem~D, whose entropy is at most $C_X\log_2(1/\varepsilon)$ for small $\varepsilon$. Local isometries preserve this entropy. A uniform extraction guarantee $\eta(\varepsilon)$ would imply $1/(512\eta(\varepsilon)^2)\le C_X\log_2(1/\varepsilon)$, proving the modulus bound. If the modulus does not tend to zero the same lower asymptotic restriction is immediate.
\end{proof}
